\documentclass[12pt,a4paper]{article}

\usepackage[utf8]{inputenc}
\usepackage[T1]{fontenc}
\usepackage[french]{babel}
\usepackage{geometry}
\usepackage{graphicx}
\usepackage{listings}
\usepackage{xcolor}
\usepackage{booktabs}
\usepackage{array}
\usepackage{hyperref}
\usepackage{titlesec}
\usepackage{fancyhdr}
\usepackage{enumitem}

\geometry{margin=2.5cm}

% ── Code style ──────────────────────────────────────────────────────────────
\definecolor{codebg}{RGB}{245,245,245}
\definecolor{codeframe}{RGB}{200,200,200}
\definecolor{keyword}{RGB}{0,0,180}
\definecolor{comment}{RGB}{80,120,80}
\definecolor{string}{RGB}{160,0,0}

\lstset{
  language=Java,
  backgroundcolor=\color{codebg},
  frame=single,
  rulecolor=\color{codeframe},
  basicstyle=\ttfamily\small,
  keywordstyle=\color{keyword}\bfseries,
  commentstyle=\color{comment}\itshape,
  stringstyle=\color{string},
  numbers=left,
  numberstyle=\tiny\color{gray},
  numbersep=8pt,
  breaklines=true,
  showstringspaces=false,
  tabsize=4,
}

% ── Header / Footer ──────────────────────────────────────────────────────────
\pagestyle{fancy}
\fancyhf{}
\rhead{Architectures Logicielles -- TP 2}
\lhead{SIQ 2025/2026}
\cfoot{\thepage}

% ── Title ────────────────────────────────────────────────────────────────────
\title{
  \vspace{-1cm}
  \textbf{TD/TP 2 -- Activité 2} \\[0.4em]
  \large Style Pipe \& Filtre \\[0.2em]
  \normalsize Architectures Logicielles -- SIQ1
}
\author{HIMEUR Idris-- DIAR Adem}
\date{2025/2026}

% ─────────────────────────────────────────────────────────────────────────────
\begin{document}

\maketitle
\tableofcontents
\newpage

% =============================================================================
\section{Introduction}

L'objectif de cette activité est de transformer l'application calculatrice développée
lors de l'Activité~1 en une architecture à base du style \textbf{Pipe \& Filtre}.
Le code de calcul, de persistance et d'interface graphique est conservé à l'identique ;
seule la \emph{structure de communication} entre les composants change.

% =============================================================================
\section{Description de l'architecture réalisée}

\subsection{Principe général}

Dans le style Pipe \& Filtre, les composants (\emph{filtres}) sont des unités de
traitement indépendantes qui ne communiquent qu'à travers des canaux de données
(\emph{pipes}). Chaque filtre :
\begin{itemize}
  \item lit ses données d'entrée depuis son \texttt{inPipe},
  \item effectue son traitement,
  \item publie son résultat dans son \texttt{outPipe}.
\end{itemize}

Aucun filtre ne connaît l'identité de ses voisins ; il ne voit que ses deux pipes.

\subsection{Les composants}

\begin{center}
\begin{tabular}{>{\bfseries}p{3.5cm} p{9.5cm}}
\toprule
Composant & Responsabilités \\
\midrule
GuiSourceFilter &
  Lit les saisies de l'utilisateur (opérandes, opération), \\
& construit un \texttt{CalculationRequest} et le pousse dans \texttt{requestPipe}. \\
\midrule
ComputeFilter &
  Calcul de la somme, du produit, du factoriel, \\
& génération d'un nombre aléatoire. \\
\midrule
TraceFilter &
  Sauvegarde de l'opération dans un fichier texte, \\
& chargement et restitution de la trace persistée. \\
\midrule
GuiOutputAdapter &
  Adaptateur frontière : tire le \texttt{TraceRecord} final du \texttt{tracePipe} \\
& et met à jour l'interface graphique. \\
\bottomrule
\end{tabular}
\end{center}

\subsection{Les pipes et les messages}

Trois pipes relient les composants. Chaque pipe est typé selon le message qu'il transporte :

\begin{center}
\begin{tabular}{>{\ttfamily}p{4.5cm} p{3.8cm} p{3cm} p{2.5cm}}
\toprule
\normalfont\bfseries Pipe & Type de message & Producteur & Consommateur \\
\midrule
requestPipe  & CalculationRequest  & GuiSourceFilter  & ComputeFilter \\
resultPipe   & CalculationResult   & ComputeFilter    & TraceFilter   \\
tracePipe    & TraceRecord         & TraceFilter      & GuiOutputAdapter \\
\bottomrule
\end{tabular}
\end{center}

\subsection{Flux d'exécution}

Lorsque l'utilisateur clique sur un bouton, \textbf{GuiSourceFilter} est déclenché :
il lit les champs de saisie, pousse un \texttt{CalculationRequest} dans \texttt{requestPipe},
puis appelle \texttt{pipeline.run()} qui va lancé par la suite \textbf{ComputeFilter}, \textbf{TraceFilter} et \textbf{GuiOutputAdapter}.

\textbf{ComputeFilter} tire la requête, calcule le résultat et pousse un
\texttt{CalculationResult} dans \texttt{resultPipe}.
\textbf{TraceFilter} tire ce résultat, le persiste et pousse un \texttt{TraceRecord}
dans \texttt{tracePipe}.
Enfin, \textbf{GuiOutputAdapter} tire le \texttt{TraceRecord} et met à jour l'affichage.

\subsection*{\colorbox{yellow!30}{\textbf{Note importante -- Composants frontières et couplage GUI}}}
\addcontentsline{toc}{subsection}{Note importante -- Composants frontières et couplage GUI}

\noindent Le style Pipe \& Filtre est conçu pour le traitement de flux de données :
chaque filtre est une unité de transformation indépendante, sans connaissance du monde extérieur.
Les données circulent dans un seul sens : elles entrent par un filtre source et sortent par un filtre sink.

\medskip
\noindent Dans une application à interface graphique, certaines opérations
(la lecture des champs de saisie et mise à jour de l'affichage) ne correspondent
pas à des transformations de données mais à des interactions avec un système externe.
Dans ce type d'architecture, ces interactions sont généralement placées
dans des \textbf{composants frontières}.

\medskip
\noindent Dans notre pipeline, ces responsabilités sont réparties entre deux filtres frontières :

\begin{itemize}
  \item \textbf{GuiSourceFilter} : filtre d'entrée qui lit les données saisies
        par l'utilisateur dans \texttt{MainGUI} et les injecte dans le pipeline via
        \texttt{requestPipe}.
  \item \textbf{GuiOutputAdapter} : filtre de sortie qui récupère le résultat final
        depuis \texttt{tracePipe} et met à jour l'interface graphique.
\end{itemize}

\noindent Ces filtres servent de frontière entre l'interface graphique et le
pipeline de traitement, tandis que les filtres internes (\texttt{ComputeFilter} et
\texttt{TraceFilter}) restent totalement indépendants de la GUI.

\medskip
\noindent La chaîne de traitement est \textbf{strictement acyclique} :
\texttt{GuiSourceFilter} ne reçoit rien du pipeline et
\texttt{GuiOutputAdapter} n'y envoie rien. Le flux de données suit un sens unique :

\begin{center}
\texttt{GuiSourceFilter} 
$\xrightarrow{\texttt{requestPipe}}$ \texttt{ComputeFilter}
$\xrightarrow{\texttt{resultPipe}}$ \texttt{TraceFilter}
$\xrightarrow{\texttt{tracePipe}}$ \texttt{GuiOutputAdapter}
\end{center}

\noindent L'orchestration est confiée à \texttt{CalculatorPipeline}, responsable
du câblage des pipes et de l'initialisation du flux de traitement. L'exécution
est déclenchée par l'injection d'un message dans le pipeline via
\texttt{GuiSourceFilter}. Chaque filtre lit les données depuis son
\texttt{inPipe} et publie son résultat dans son \texttt{outPipe}. Aucun filtre
ne connaît ses voisins ; seul le coordinateur définit la structure globale du pipeline.



% =============================================================================
\section{Style d'échange utilisé}

Le style d'échange utilisé est le \textbf{PIPE} : les données circulent sous
forme de \textbf{flux unidirectionnel} entre les composants.
Chaque producteur \textbf{pousse} (\texttt{push}) son message dans le pipe ;
le consommateur suivant le \textbf{tire} (\texttt{pull}) pour le traiter.
Les filtres ne se connaissent pas directement -- seuls les pipes les relient.

% =============================================================================
\section{Vues de l'architecture}

\subsection{Vue de développement (vue des modules)}

\begin{figure}[h!]
  \centering
  \includegraphics[width=1\textwidth]{calculator_architecture.png}
  \caption{Vue de développement}
  \label{fig:vue_developpement}
\end{figure}

\subsection{Vue C\&C (Composants et Connecteurs)}

\begin{figure}[h!]
  \centering
  \includegraphics[width=1\textwidth]{calculator_pipeline.png}
  \caption{Vue C\&C -- Pipe \& Filtre}
  \label{fig:vue_cc}
\end{figure}


% =============================================================================
\section{Comparaison Activité 1 vs Activité 2}

\begin{center}
\begin{tabular}{p{7cm} p{7cm}}
\toprule
\textbf{Activité 1 -- Appel de méthode direct} &
\textbf{Activité 2 -- Pipe \& Filtre} \\
\midrule
\textbf{(+) Performance} : pas de sérialisation ni de
tampon intermédiaire ; l'appel est synchrone et instantané.
&
\textbf{(+) Indépendance des composants} : aucun filtre
ne connaît ses voisins ; on peut remplacer CMP2 sans toucher
CMP1 ni CMP3.
\\[0.5em]
\textbf{(+) Simplicité de la structure} : une seule classe
centrale relie tout.
&
\textbf{(+) Lisibilité architecturale} : la vue C\&C montre
immédiatement le flux de données et les responsabilités.
\\[0.5em]
\textbf{(-) Couplage fort} : la GUI connaît directement
\texttt{Calcul} et \texttt{Trace} ; tout changement se propage.
&
\textbf{(+) Réutilisabilité et testabilité} : chaque filtre
peut être testé isolément avec un pipe bouchon.
\\[0.5em]
\textbf{(-) Difficile à faire évoluer} : ajouter une étape
de traitement implique de modifier la GUI ou le coordinateur.
&
\textbf{(+) Extensibilité} : ajouter un filtre revient à créer
une nouvelle classe et à l'insérer dans la chaîne de pipes.
\\[0.5em]
&
\textbf{(-) Transformations interactives difficiles} :
le retour du résultat vers la GUI nécessite un pipe
supplémentaire (\texttt{tracePipe} $\rightarrow$ GuiFilter),
ce qui complexifie légèrement le câblage.
\\[0.5em]
&
\textbf{(-) Overhead léger} : création d'objets messages
à chaque opération (minime en pratique).
\\
\bottomrule
\end{tabular}
\end{center}

\medskip
\noindent\textbf{Conclusion.}
L'Activité~1 est plus performante et plus simple pour une petite application.
L'Activité~2 est architecturalement plus robuste : les composants sont vraiment
indépendants, le flux de données est explicite, et l'application peut évoluer
sans modifier les composants existants.

% =============================================================================
\end{document}